\documentclass[
    a4paper, % Paper size, use either a4paper or letterpaper
    10pt, % Default font size, can also use 11pt or 12pt, although this is not recommended
    unnumberedsections, % Comment to enable section numbering
]{LTJournalArticle}

\usepackage{subfig} % grace added this to add plots in side by side! sorry didn't know where else to put it. feel free to move / remove
\usepackage{amsmath}
\usepackage{cleveref}
\usepackage{stfloats}
\usepackage{float}

\hypersetup{
    colorlinks=true,
    linkcolor=magenta,
    citecolor=magenta,
    urlcolor=magenta
}

%\AtEveryBibitem{\clearfield{title}}

\runninghead{The Colour of \textsc{simba}'s Universe} 

\begin{document}
\epsfxsize=40mm                         % Size of crest
\thispagestyle{empty} % This removes the running head from the top corner
\begin{minipage}[b]{110mm}
        {\Huge\bf School of Physics \\and Astronomy
        \vspace*{17mm}}
\end{minipage}
\hfill
\begin{minipage}[t]{40mm}               
        \makebox[40mm]{
        
\includegraphics[width=40mm]{figures/manuscript/UoEcrest.pdf}}
\end{minipage}
\par\noindent                                           % Centre Title, and name
\vspace*{2cm}
\begin{center}
        \huge\bf The Colour of \textsc{\textmd{simba's}} Universe: The Extragalactic Background Light Prediction of the \textsc{\textmd{simba}} Simulations
\end{center}
\vspace*{1cm}
\begin{center}
        \bf 
        \large
        Siddhant Pujni
        \vspace{0.5cm} \\
        March 2026                         % Submission Date
\end{center}
\vspace*{5mm}


\begin{abstract}

    The extragalactic background light traces the integrated emission from all galaxies over cosmic time and provides a key constraint on models of galaxy formation. In this work, the EBL is predicted using the \textsc{simba} cosmological simulations by constructing a light cone over $0 \leq z \leq 7$ and modelling optical/near infrared, far infrared, and radio emission components. The resulting energy distribution reproduces the expected bimodal structure, with a stellar peak in the optical and a dust reprocessed peak in the far infrared. While the overall shape agrees with observations, the optical emission is under predicted and the radio component overestimated, whereas the far infrared shows reasonable agreement. The majority of the EBL is found to originate from $z < 3$. These results highlight the ability of simulations to capture the global energy budget, while emphasising the importance of improved modelling of dust and feedback processes.

\end{abstract}

\vspace*{4cm}

\vfill
\begin{center}
    {\bf Supervisor:} Professor Romeel Davé               
\end{center}
\newpage

\twocolumn
\setcounter{page}{1} % Set this page as the first numbered page


\section{Introduction}

%In modern cosmology, two primary approaches are employed: hydrodynamical simulations and semi-analytic models (SAMs), which differ in both their computational complexity and the level of physical detail used to describe galaxy formation and evolution \parencite{neisteinHydrodynamicalSimulationsSemianalytic201}.

\textit{Olbers' paradox} was a fundamental question in cosmology, it questioned why the night sky appears dark despite the presence of an effectively infinite number of luminous sources in an 'infinite' universe. The resolution lies in the finite age of the universe and the effects of cosmic expansion. Photons emitted by distant sources are redshifted as the universe expands, shifting their energy to longer wavelengths and reducing their observed intensity. Consequently, the sky is not truly dark but instead filled with radiation across the entire electromagnetic spectrum. This diffuse radiation field is known as the extragalactic background light (EBL), representing the integrated emission from all sources over cosmic time. It is typically expressed as a specific intensity in units of $\mathrm{nW\ m^{-2}\ sr^{-1}}$ and has been measured across multiple wavelength regimes in several key studies \parencite{longairUniversalElectromagneticBackground1972,kashlinskyCosmicInfraredBackground2005,wessonOlberssParadoxSpectral1991,hauserCosmicInfraredBackground2001a,bernsteinOpticalExtragalacticBackground2007,vavrycukUniverseOpacityEBL2017}, however, the data sets have occasionally had their reliability questioned \parencite{mattila1mmDiscontinuityExtragalactic2006}.

A closely related quantity is the cosmic spectral energy distribution (CSED), which describes the total electromagnetic power produced per unit comoving volume as a function of wavelength at a given epoch. While the EBL represents the time integrated light from all redshifts, the CSED provides a snapshot of the energy output of the universe at a specific redshift. In this sense, the EBL can be interpreted as the integral of the evolving CSED over cosmic time,

\begin{equation}
    \mathrm{EBL} = \int \mathrm{CSED}(z)\frac{dt}{dz}dz.
\end{equation}

Physically, the CSED arises from the combined spectral energy distributions (SEDs) of all galaxies within a representative cosmological volume, encoding information about their star formation history, dust content, and active galactic nucleus (AGN) activity \parencite{andrewsModellingCosmicSpectral2018}.

One may measure the EBL by dividing it into its principal electromagnetic components; gamma-rays, X-rays, ultraviolet, optical, infrared, and radio, which enables a more detailed interpretation of the physical processes contributing to the total cosmic emission \parencite{hillSpectrumUniverse2018}. Considerable observational effort has been devoted into measurements of the CSED to reconstruct the total energy output of galaxies in the ultraviolet \parencite{wilsonStarFormationHistory2002,budavariUltravioletLuminosityFunction2005,wyderUltravioletGalaxyLuminosity2005,robothamGALEXSDSSNUV2011}, optical \parencite{norberg2dFGalaxyRedshift2002,blantonGalaxyLuminosityFunction2003,montero-dortaSDSSDR6Luminosity2009}, and infrared \parencite{kochanekKBandGalaxyLuminosity2001,cole2dFGalaxyRedshift2001,saunders60mFarinfraredLuminosity1990,marchettiHerMESSubmillimetreLocal2016}. However, combining measurements from independent studies introduces systematic inconsistencies arising from differences in source selection, survey depth, and methodology, compounded by the effects of cosmic variance across different sight lines \parencite{driverGalaxyMassAssembly2012,crossBivariateBrightnessFunction2002}. A more rigorous approach is to construct the CSED by combining measurements from large, multiwavelength datasets that span a wide range of cosmic volumes and redshifts, as demonstrated by a number of wide field surveys \parencite{hillUgrizYJHKLuminosityDistributions2010,driverGalaxyMassAssembly2011,andrewsModellingCosmicSpectral2018}. 

The shape and normalisation of the CSED, as well as its evolution with redshift, provide strong observational constraints on models of galaxy formation and evolution by encoding the cumulative history of star formation, dust reprocessing, and nuclear activity across the observable universe \parencite{dominguezExtragalacticBackgroundLight2011}. The integrated EBL also provides a probe of the global star formation history \parencite{fardalEvolutionaryHistoryStars2007}, as the total photon budget must be consistent with the number and mass of stars formed over cosmic time. The EBL is expected to exhibit two dominant peaks. The first lies in the optical/near infrared, arising from the direct emission of stellar populations, particularly young, massive stars that radiate strongly at short wavelengths \parencite{schlegelMapsDustInfrared1998}. A second peak in the far infrared results from interstellar dust grains absorbing ultraviolet and optical radiation and re-emitting this energy thermally \parencite{calzettiDustContentOpacity2000}. 

Overall, the different wavelength regimes of the EBL trace distinct physical processes, each contributing a characteristic signature. Ultraviolet and optical emission is dominated by stellar photospheres, with young, massive stars contributing strongly at short wavelengths, making this regime a probe of recent star formation \parencite{kennicuttStarFormationMilky2012}. At longer wavelengths, far infrared emission arises from interstellar dust that absorbs stellar radiation and re-emits it thermally. The mid infrared is characterised by emission from polycyclic aromatic hydrocarbons (PAHs), complex carbon based molecules that are stochastically heated by ultraviolet photons and radiate through vibrational transitions \parencite{peetersPolycyclicAromaticHydrocarbon2023}. Radio emission arises through two distinct mechanisms: star formation and AGN activity. Radio observations trace the star formation within galaxies due to the synchrotron radiation from cosmic ray electrons accelerated in supernova remnants and free-free (bremsstrahlung) emission from ionised gas in HII regions indicative of young hot stars \parencite{kopsacheiliNewXraySupernova2025}. Additionally, radio loud AGN inject energy via relativistic jets that dominate the bright end of the radio source population \parencite{heckmanCoevolutionGalaxiesSupermassive2014}.

The EBL can also be interpreted in terms of the dominant photon production pathways operating after recombination, following the decoupling of the cosmic microwave background (CMB). These processes include: (i) star formation, which drives emission across the UV, optical, and radio regimes through stellar radiation and associated particle processes \parencite{partridgeAreYoungGalaxies1967, tinsleyGalaxyCountsColorredshift1977}; (ii) accretion onto supermassive black holes, which powers AGN and contributes significantly at high energies \parencite{silvaConnectingCosmicInfrared2004, ballantyneContributionActiveGalactic2007}; (iii) AGN jets, producing extended radio emission via synchrotron radiation \parencite{khaireNewSynthesisModels2019}; and (iv) dust heated by both star formation and AGN activity, which re-emits absorbed energy in the infrared \parencite{lagacheDustyInfraredGalaxies2005}. Each of these pathways imprints a distinct spectral signature on the EBL, such that a well modelled multi wavelength prediction provides a powerful constraint on the star formation history, dust content, and black hole accretion processes across cosmic time.

Existing predictions of the EBL from hydrodynamical simulations remain relatively limited in scope, with \cite{baesCosmicSpectralEnergy2019} providing one of the few examples, restricted to the local universe. While EBL predictions from semi-analytic models, such as \cite{gilmoreSemianalyticModelingEBL2012}, are more widely available, these approaches typically rely on simplified prescriptions for galaxy formation and do not fully capture the self-consistent interplay between feedback processes, black hole growth, and interstellar medium enrichment inherent to hydrodynamical simulations.

The present work addresses these limitations by constructing a multi wavelength EBL prediction using the \textsc{simba} cosmological simulations, spanning the optical, far infrared, and radio regimes over a wide redshift range. A light cone approach is employed to recreate the geometry of real observations, representing a novel extension of hydrodynamical simulation based EBL studies. The wavelength ranges are selected to capture the dominant contributions from stellar emission, dust reprocessing, and radio activity, enabling a physically motivated reconstruction of the global EBL.

The primary aim of this project is to construct a prediction for the EBL from the \textsc{simba} simulations and compare to observational estimates from Driver et al, an upcoming paper that presents a heterogeneous mix of data from the literature. To achieve this, galaxy catalogues derived from \textsc{simba} are used to compute the contributions to the cosmic SED across different wavelength regimes. Light cone techniques are employed to account for the evolution of the galaxy population with redshift and to produce a realistic observational analogue. The resulting SED is then analysed as a function of wavelength and redshift, and statistical uncertainties are estimated using jackknife resampling of the light cone. By comparing the simulated EBL to observational constraints, this work aims to assess the ability of the SIMBA simulation to reproduce the global energy output of the universe and to identify potential discrepancies that may inform future improvements in galaxy formation models.

\section{Methodologies}

\subsection{The \textsc{\textmd{simba}} Simulations}

Hydrodynamical simulations serve as theoretical laboratories for exploring the formation and evolution of structure in the Universe. By modelling the complex interplay between gravity, gas dynamics, star formation, and feedback processes, these simulations provide a controlled framework within which galaxy evolution can be studied. As observational capabilities through the James Webb Space Telescope (JWST), \textit{Euclid}, and the Dark Energy Spectroscopic Instrument (DESI), simulations play an increasingly important role in interpreting observations and refining our understanding of the underlying astrophysical processes. The \textsc{simba} simulations \parencite{daveSimbaCosmologicalSimulations2019} are hydrodynamical state-of-the-art simulations of galaxy formation, built on \textsc{gizmo} \parencite{hopkinsNewClassAccurate2015} using its meshless finite mass (MFM) method. Building on the previous generation \textsc{mufasa} simulations with improvements to the subgrid prescriptions for both star formation and active galactic nucleus (AGN) feedback via modified code from \textsc{gadget-3} \parencite{springelCosmologicalSimulationCode2005} it evolves dark matter and baryonic elements together including gravity and pressure forces.

\textsc{simba} approaches star formation based on molecular gas fraction following the $\mathrm{H_2}$ Schmidt-Kennicutt relation with a fixed star formation efficiency $\epsilon_\star=0.02$ \parencite{robertc.kennicuttGlobalSchmidtLaw1998}. The star formation rate (SFR) is based on the $H2$ fraction which is computed via the \cite{krumholzCOMPARISONMETHODSDETERMINING2011} model. Stellar feedback is present via metal loaded kinetic winds that enrich the interstellar medium (ISM) for future star formation. Both torque limited \parencite{angles-alcazarGravitationalTorquedrivenBlack2017} and Bondi accretion \parencite{bondiSphericallySymmetricalAccretion1952,bondiMechanismAccretionStars1944} are used to model super massive black hole (SMBH) growth based on a cold accretion mode or a hot halo mode respectively. The resulting energetic feedback is modelled kinetically depending on the Eddington ratio which serves to quench galaxies along with X-ray energy feedback based on \cite{choiRADIATIVEMOMENTUMBASEDMECHANICAL2012}.

Finally, \textsc{simba} uses a self-consistent dust framework evolving directly from initial conditions, which is responsible for production, growth, and destruction \parencite{liDusttoGasDusttoMetalsRatio2019}. All dust grains are assumed to have the same constant size of $0.1 \mu m$ and have their motions coupled to gas flows. Dust production is informed by Type II supernovae (SNe) and Asymptotic Giant Branch stars, and further growth via condensation from metals, while its destruction can occur from sputtering, consumption by star formation, or SNe shocks.

The primary \textsc{simba} run is used for this work containing $1024^3$ dark matter and $1024^3$ baryonic particles in a comoving volume $100h^{-1}$Mpc in a redshift range $0<z<249$. The resulting mass evolution is $9.6 \times 10^7 M_\odot$ for dark matter particles and $1.82 \times 10^7 M_\odot$ for baryonic gas elements. The simulations follow a $\Lambda$CDM cosmology consistent with Planck16 \parencite{planckcollaborationPlanck2015Results2016}, and parameters $\Omega_m = 0.3$, $\Omega_\Lambda = 0.7$, $\Omega_b = 0.048$, and $H_0 = 68 \mathrm{kms^{-1}Mpc^{-1}}$. Simulation runs output 151 snapshots from $z=20 \rightarrow 0$ which are processed into \textsc{caesar} catalogues which contain the necessary pre-computed properties for galaxies and halos.\textsc{ caesar}\footnote{caesar.readthedocs.io/en/latest/index.html} is a python based \textsc{yt} extension for analysing outputs from cosmological simulations through an intuitive interface without the original simulation binary. It outputs a \textsc{hdf5} catalogue that contains the galaxy and halo properties used for the analysis in this work. One may also conduct photometry for galaxy spectra using \textsc{fsps}.

\subsection{Light Cones}
\label{subsec:lcConstruction}

A light cone is a geometric construct typically seen in special relativity that represents the path of light travelling through spacetime from distant sources to an observer. In our context of galaxy surveys and cosmological simulations, a light cone catalogue mimics what a real telescope would observe: galaxies at different distances are seen at different cosmic times, so we account for the fact that the galaxy population evolves with cosmic time and that more distant galaxies contribute less flux due to the inverse square law and cosmological dimming. Constructing a light cone with the \textsc{simba} simulations involves stitching together periodic snapshots, each representing the universe at a specific redshift, we adopt light cone generation as discussed in \cite{lovellReproducingSubmillimetreGalaxy2021}\footnote{github.com/christopherlovell/simba\_lightcone\_generation}. The resulting light cone allows computation of the respective SED for each electromagnetic spectrum component, hence we obtain a prediction for the total EBL that would be measured by an observer.

\begin{figure*}
    \centering
    \includegraphics[width=0.9\linewidth]{figures/lightcone/lightcone_wedge_three_panel.png}
    \caption{Wedge diagrams of the \textsc{simba} light cone showing comoving radial distance (Mpc) on the horizontal axis and comoving transverse distance (Mpc) on the vertical axis, with a secondary axis along the top indicating the corresponding redshift. Each point represents an individual galaxy, with $240{,}883$ galaxies in total. Panel (a) is coloured by $V$-band apparent magnitude $m_V$; panel (b) by far IR luminosity $L_\mathrm{FIR}$; and panel (c) by 1.4\,GHz radio flux density $S_\mathrm{1.4\,GHz}$.}
    \label{fig:lcWedge}
\end{figure*}

\Cref{fig:lcWedge} shows wedge diagrams of the \textsc{simba} light cone for the three spectral components considered in this work. The expanding geometry with radial distance reflects the growth of the comoving volume per redshift interval, naturally producing a higher surface density of galaxies at larger distances. The discrete vertical banding visible across all three panels is a construction artefact, arising from the finite number of simulation snapshots tiled together to build the light cone. In all three panels, the galaxy distribution is clearly concentrated at $z < 3$, with the number density dropping off significantly beyond this.

Panel (a) shows that optically bright galaxies, with the smallest $m_V$, are predominantly local, as expected given the dimming of apparent magnitude with distance. Panel (b) reveals that the most infrared luminous galaxies are spread across a broader redshift range, peaking around $z \sim 1-3$. Panel (c) shows that the brightest radio sources are likewise concentrated at low redshift, reflecting the dominance of nearby radio loud AGN in the radio background.

With the light cone constructed, it serves as the geometric and observational framework through which the EBL is computed. For each spectral component, the relevant physical properties of every galaxy in the light cone are used to compute a flux contribution, with the light cone geometry naturally encoding the effects of luminosity distance, cosmological dimming, and the $(1+z)$ redshifting of frequencies between the rest and observed frames. These individual galaxy contributions are then summed and divided by the solid angle of the light cone to produce a sky averaged specific intensity, yielding the predicted EBL as a function of observed wavelength. The methodology for computing each spectral component from these galaxy properties is described in the sections that follow.

\subsection{The Optical/NIR Regime}

Using the light cone construction described above we estimate the optical/near infrared (optical/NIR) cosmic background by integrating the observed flux from all simulated galaxies. Using effective wavelengths for discrete photometric bands in \textsc{fsps}, a sky averaged intensity $I_\nu$ is calculated as a function of observed wavelength. We plot the apparent magnitude including dust attenuation computed via \textsc{caesar}. We use the apparent magnitude as they encode the flux dimming with distance and are a natural choice for the observed intensity EBL.

The optical/NIR range itself is determined operationally by the set of available filters. Each filter is mapped to an effective wavelength $\lambda_{\mathrm{eff}}$ using the \textsc{fsps} database. This is required so that band integrated photometric points are on a common wavelength axis for unit conversion. 

After $\lambda_{\mathrm{eff}}$ were determined the next step was to convert the apparent magnitudes stored in the \textsc{caesar} catalogues into a flux density. For each galaxy in the light cone, for each filter, the apparent magnitude, $m$, is converted to a flux density in the AB system via:

\begin{equation}
    f_\nu({\mathrm Jy}) = 3631 \times 10^{-m/2.5}.
\end{equation}
\noindent Next, for every galaxy appearing in the light cone, their flux densities are summed in each filter to get an integrated flux density, $F_\nu$. This accounts for any luminosity distance or K-corrections as we are using apparent magnitudes. We then further convert this to a mean specific intensity by dividing $F_\nu$ by the light cone solid angle,

\begin{equation}
    \Omega\ ({\mathrm sr}) = A\left(\frac{\pi}{180}\right)^2,
    \label{eq:solidAngle}
\end{equation}

\noindent where $A$ is sky area of the light cone. This yields $I_\nu$ presented in units of $\mathrm{erg\ s^{-1}\ cm^{-2}\ Hz^{-1}}$ for each filter $\lambda_\mathrm{eff}$. Finally, $\nu I_\nu$ is calculated using $\nu =\frac{c}{\lambda_{\mathrm{eff}}}$ to show how much energy is contained within a given range of frequencies. It is particularly useful because it allows identification of which frequency band dominates the total energy emission (the peak of the plot represents the maximum power). A dust free optical curve was also produced by removing the attenuation correction, allowing the effect of dust reprocessing on the UV/optical emission to be isolated. Since \textsc{caesar} stores no dust apparent magnitudes directly, the correction was applied the same as described above.

\subsection{The Far Infrared Regime}

The far infrared (far IR) cosmic background is estimated by integrating the redshifted thermal dust emission from simulated galaxies along the constructed light cone. A full scale radiative transfer approach was not possible due to the scope of the project and time constraints, hence a simplified approach is used. Dust emission deviates from a perfect blackbody, being cooler and less efficient at shorter wavelengths, and so each galaxy's far IR SED is modelled as a single temperature modified blackbody (MBB), or greybody, in which the Planck function $B_\nu(T)$ is scaled by a frequency dependent emissivity power law. Since the absorption cross section of dust grains drops dramatically at long wavelengths this formulation assumes the dust is optically thin ($\tau \ll 1$), meaning radiation passes through without significant absorption \parencite{hildebrandDeterminationCloudMasses1983}.

Following \cite{caseyFarinfraredSpectralEnergy2012}, but in wavelength space via $\lambda = c/\nu$ (which introduces a factor of $\nu^2$ from the Jacobian), the rest frame optically thin MBB is

\begin{equation}
    S(\lambda) = N \cdot \left(\frac{100\,\mu\mathrm{m}}{\lambda}\right)^{\!\beta} 
    \cdot \frac{2hc^2}{\lambda^5} \cdot \frac{1}{\exp\!\left(\dfrac{hc}{\lambda k_B T}
    \right) - 1},
    \label{eq:dustmbb}
\end{equation}

\noindent where $N$ is a normalisation constant, $\left(100\,\mu\mathrm{m}/\lambda\right)^\beta$ is the dust emissivity power law anchored at $100\,\mu\mathrm{m}$, and $\beta$ is the spectral emissivity index (how efficiently a surface emits radiation at a specific wavelength). The normalisation $N$ is set such that the integral of $S(\lambda)$ over wavelength equals the galaxy's bolometric far IR luminosity $L_\mathrm{FIR}$ (the total energy absorbed by dust extinction) computed by \textsc{caesar}. This ensures that more luminous galaxies contribute proportionally more flux to the background, while $T$ and $\beta$ govern the spectral distribution of that luminosity \parencite{caseyFarinfraredSpectralEnergy2012}.

The equivalent dust temperature, $T_\mathrm{eqv}$, is estimated using the relation presented in \cite{liangDustTemperaturesHighredshift2019}, which was calibrated against the \textsc{FIRE} simulations. These simulations are known for high resolution and their ability to accurately model the internal structure of galaxies and the ISM \parencite{feldmannAreStarFormation2017}. Given the similar level of physical detail in the treatment of the ISM within \textsc{simba}, this relation provides a suitable prescription for estimating dust temperatures in our analysis. Accordingly, a dust temperature is assigned to each galaxy as

\begin{equation}
    \log\left(\frac{T_\text{eqv}}{25\,\text{K}}\right) = a_{\mathrm{dust}} + 
    b\log\!\left(\frac{\delta_\text{dzr}}{0.4}\right) + c\log(1+z),
    \label{eq:equivalentDustTemp}
\end{equation}

\noindent where $b$ and $c$ are the best fitting regression parameters from that work, derived for an optically thin MBB evaluated at $870\,\mu\mathrm{m}$ at a fixed $\beta = 2$ which are applied for this analysis as well. $T_\mathrm{eqv}$ is not a resolved temperature distribution but a single effective temperature that, when paired with $L_\mathrm{FIR}$, reproduces the galaxy's far IR emission. The dust-to-metal ratio, $\delta_\mathrm{dzr}$, is computed per galaxy as

\begin{equation*}
    \delta_{\mathrm{dzr}} = \frac{M_\mathrm{dust}}{Z \cdot M_{\mathrm{gas}}},
\end{equation*}

\noindent where $M_\mathrm{dust}$ dust mass of the galaxy, $Z$ is its metallicity, and $M_\mathrm{gas}$ is the gas mass, obtained from the \textsc{caesar} catalogues \parencite{liDusttoGasDusttoMetalsRatio2019}. The 
overall temperature scaling parameter $a_{\mathrm{dust}}$ is treated as a free parameter and determined by fitting the simulated SED to observational data. The far IR spectra are restricted to $50$-$3000\,\mu\mathrm{m}$ to exclude contamination from the optical/NIR and radio regimes. Model predictions are interpolated onto the observed wavelength grid and $a_{\mathrm{dust}}$ is determined by minimising the $\chi^2$ statistic,

\begin{equation}
    \chi^2(a_{\mathrm{dust}})=\sum_i \left[\frac{M_i(a_{\mathrm{dust}})-D_i}
    {\sigma_i}\right]^2,
    \label{eq:continiousObjectiveFunction}
\end{equation}

\noindent where $D_i$ and $\sigma_i$ are the observed value and uncertainty at wavelength $\lambda_i$, and $M_i(a_{\mathrm{dust}})$ is the model prediction \parencite{jeyakumarContinuousOptimizationCurrent2006}. Since $a_{\mathrm{dust}}$ is sampled on a discrete grid, a secondary interpolation is performed along $a_{\mathrm{dust}}$ to allow a bounded one dimensional minimiser to find the best fitting value continuously.

\begin{figure*}
    \centering
    \includegraphics[width=0.8\linewidth]{figures/farIR/farIR_bg_m100n1024_diagnostic.png}
    \caption{Top: far IR SEDs produced by varying $a_\mathrm{dust}$, with the best fit value $a_\mathrm{dust} = -0.0173(6)$ shown alongside curves spanning $\pm50\sigma$. Observational EBL data are overplotted for reference. Bottom: the corresponding effective dust temperature distributions for each $a_\mathrm{dust}$ value. A cosmic dust temperature of $44.7 \pm 7.7$\,K is found. }
    \label{fig:tempDiagnostic}
\end{figure*}

\Cref{fig:tempDiagnostic} shows diagnostic plots for the dust temperature scaling parameter $a_\mathrm{dust}$, illustrating how different values of $a_\mathrm{dust}$ map to different effective temperature distributions. To show a meaningful spread around the best fit solution, curves spanning $\pm50\sigma$ are plotted alongside the observational EBL data.

The resulting temperature distribution reflects the redshift evolution of the dust obscured galaxy population. In the local universe, mean dust temperatures are expected to be low ($\sim20-30$\,K), consistent with quiescent galaxies such as the Milky Way. At higher redshifts, approaching cosmic noon ($z \sim 2$), the galaxy population is on average more actively star forming, with younger, hotter stellar populations driving elevated dust temperatures \parencite{parenteDriversCosmicDust2026}. This redshift dependence is encoded in $a_\mathrm{dust}$ and is reflected in the median temperature of the distribution being higher than the local value, while the low temperature tail of the histogram confirms that cooler, locally representative dust temperatures are still well sampled.

To account for redshift, the normalised SED is evaluated at the rest frame wavelength $\lambda_\mathrm{rest} = \lambda_\mathrm{obs}/(1+z)$, and the observed flux density is

\begin{equation}
    f_\lambda(\lambda_{\mathrm{obs}}) =
    \frac{L_\lambda(\lambda_{\mathrm{rest}})}{4\pi d_L(z)^2 \cdot (1+z)},
    \label{eq:luminosityDistanceRelation}
\end{equation}

\noindent where $d_L(z)$ is the luminosity distance and the factor of $(1+z)$ accounts for the stretching of wavelength bins between the rest and observed frames \parencite{hoggDistanceMeasuresCosmology2000,gabrielliStatisticalPhysicsCosmic2005}. The sky averaged intensity is then obtained by dividing the total flux by the light cone solid angle in \cref{eq:solidAngle}. A logarithmic wavelength grid is adopted throughout to avoid undersampling at long wavelengths and oversampling at short wavelengths.

\subsection{The Radio Regime}

For the radio emission background we consider a broad interval from 10 MHz to 100 GHz to account for contributions from both SFR and AGN activity. Firstly, star formation observations are decomposed into two components as discussed earlier: non thermal synchrotron emission from cosmic ray electrons accelerated in supernova remnants, and the thermal bremsstrahlung emission \cite{condonRadioEmissionNormal1992}. The total radio luminosity is modelled following \cite{thomasRadioGalaxyPopulation2021} as $P_\nu = P_\mathrm{non-thermal} + P_\mathrm{thermal}$, where

\begin{multline}
    P_{\mathrm{non\text{-}thermal}} (\mathrm{W\,Hz^{-1}})
= 5.3 \times 10^{21}
\left(\frac{\nu}{\mathrm{GHz}}\right)^{-0.8} \\
\times \mathrm{SFR}[M \geq 5\,M_\odot],
\label{eq:nonThermalRadio}
\end{multline}

\begin{multline}
P_{\mathrm{thermal}} (\mathrm{W\,Hz^{-1}})
= 5.5 \times 10^{20}
\left(\frac{\nu}{\mathrm{GHz}}\right)^{-0.1} \\
\times \mathrm{SFR}[M \geq 5\,M_\odot].
\label{eq:thermalRadio}
\end{multline}

\noindent Both relations are expressed in terms of $\mathrm{SFR}[M \geq 5\,M_\odot]$, which is the rate of mass going into stars more massive than $5\,M_\odot$, since only these stars produce the supernovae and ionising photons that drive the radio emission. This is obtained from the total SFR by

\begin{equation}
    \mathrm{SFR}[M \geq 5\,M_\odot] = f_{M \geq 5} \times \mathrm{SFR_{total}},
\end{equation}

\noindent where $\mathrm{SFR_{total}}$ is taken directly from the \texttt{caesar} catalogues and $f_{M \geq 5}$ is the initial mass function (IMF) mass fraction above $5\,M_\odot$,

\begin{equation}
    f_{M \geq 5} = \frac{\int_{5}^{100}m\,\xi(m)\,\mathrm{d}m}{\int_{0.1}^{100}m\,\xi(m)\,\mathrm{d}m}.
    \label{eq:chabrierMassFraction}
\end{equation}

\noindent It is integrated between 0.1 and 100\,$M_\odot$, representing the full observed range of stellar masses. The lower bound reflects the minimum mass at which hydrogen burning can sustain a star, while the upper bound captures the most massive stars known to form before stellar winds and radiation pressure prevent further growth. The numerator is restricted to stars above 5\,$M_\odot$, as only these sufficiently massive stars evolve quickly enough to end their lives as core collapse supernovae, producing the synchrotron radiation and ionising photons that drive the radio emission \parencite{bietenholzRadioLuminosityrisetimeFunction2021}. $\xi(m)$ is the \cite{chabrierGalacticStellarSubstellar2003} IMF, 
which takes the form of a power law above $1\,M_\odot$,

\begin{equation}
    \xi(m) = B\,m^{-2.3}, \quad m > 1\,M_\odot,
    \label{eq:chabrierPowerLaw}
\end{equation}

\noindent and a log-normal distribution below $1\,M_\odot$,

\begin{equation}
    \xi(m) = \frac{A}{m}\exp{\left(-\frac{(\log_{10}m - \log_{10}0.08)^2}
    {2 \times 0.69^2}\right)}, 
    \quad m \leq 1\,M_\odot.
    \label{eq:chabrierLogNormal}
\end{equation}

\noindent The constants $A$ and $B$ are chosen to ensure continuity at $1\,M_\odot$ and normalisation over the full mass range. Since the log-normal form has no analytic integral, the denominator of \cref{eq:chabrierMassFraction} is evaluated numerically. The dependence on high mass stars of the radio luminosity calibration stems from their ionising photon production which drives bremsstrahlung, and them typically ending their lives as core collapse supernovae which produce the synchrotron radiation observed \parencite{shenoyWhatFactorsShape2026}.

The AGN contribution to the radio background is computed from the black hole accretion rates of individual galaxies in the \textsc{simba} simulation. For each galaxy, the bolometric radio power is derived from the black hole accretion rate as presented in \cite{kordingMeasuringAccretionRate2007},

\begin{equation}
    \frac{P_\mathrm{Rad}}{10^{30}\mathrm{erg}\text{ }\mathrm{s^{-1}}} = \left(\frac{\dot{M}_{BH}}{4\times10^{17}\mathrm{g}\text{ }\mathrm{s^{-1}}}\right)^\frac{17}{12},
    \label{agnRadioEmission}
\end{equation}

\noindent where $\dot{M}_{\mathrm{BH}}$ is the accretion rate of the most massive black hole in each galaxy as catalogued by \textsc{caesar}. Since the bolometric radio power represents the monochromatic luminosity $P_{\mathrm{Rad}} \sim \nu P_\nu$ evaluated at 1.4 GHz, the spectral luminosity is obtained by normalising as $P_\nu(1.4\;\mathrm{GHz}) = P_{\mathrm{Rad}} / \nu_{\mathrm{ref}}$, where $\nu_{\mathrm{ref}} = 1.4\;\mathrm{GHz}$. To evaluate the AGN radio emission across a frequency range, we extrapolate assuming a power law spectral energy distribution of the form

\begin{equation}
    P_\nu(\nu) = P_\nu(1.4\;\mathrm{GHz})\,\left(\frac{\nu}  {1.4\;\mathrm{GHz}}\right)^{\alpha},
\end{equation}

\noindent adopting a spectral index of $\alpha = -0.7$ characteristic of radio loud AGN jets \parencite{thomasRadioGalaxyPopulation2021,angladaRadioJetsYoung2018}. The contribution of each galaxy to the observed specific intensity is then computed by evaluating $P_\nu$ at the rest frame frequency $\nu_{\mathrm{rest}} = \nu_{\mathrm{obs}}(1+z)$ and converting to observed flux density in frequency space using \cref{eq:luminosityDistanceRelation}. The AGN flux contributions from all galaxies in the light cone are summed and divided by the survey solid angle in \cref{eq:solidAngle} to obtain the AGN component of the radio specific intensity as a function of observed frequency.

%$F_\nu = (1+z)\,P_\nu(\nu_{\mathrm{rest}}) / (4\pi d_L^2)$, where $d_L$ is the luminosity distance. 

\subsection{Uncertainty Estimation}

Sources of statistical uncertainty such as the sampling variance arising from the finite volume of the light cone and the inherent cosmic variance of the underlying matter distribution have not yet been addressed. To quantify this, we apply a jackknife resampling method \parencite{bergerJackknifeVarianceEstimator2007}, dividing the light cone into $N = 16$ spatial quadrants defined by quantile bins in right ascension and declination. For each realisation $k = 0, \ldots, 15$, the $k^{\mathrm{th}}$ quadrant is excluded and the full EBL spectrum is recomputed from the ones remaining. The jackknife variance is then estimated as:

\begin{equation}
    \sigma_\mathrm{jk}^2 = \frac{N-1}{N} \sum_{k=0}^{N-1} \left(x_k - \bar{x}\right)^2,
\end{equation}

\noindent where $x_k$ is the EBL estimate from the $k^{\mathrm{th}}$ realisation and $\bar{x}$ is the mean over all realisations. This uncertainty is computed independently at each wavelength bin and propagated through to the final SED.

There are three possible limitations to jackknife error sampling. First, the quadrant sub samples are assumed to be spatially independent, however, this assumption breaks down on scales comparable to the largest structures within the simulation volume. Second, the jackknife method captures only statistical uncertainty and does not account for systematic effects, such as those arising from the adopted dust prescription or incomplete snapshot coverage. Third, the number of realisations is inherently limited by the finite simulation volume; a larger volume would enable finer spatial subdivision and yield a more robust estimate of the variance.
 
\section{Results \& Discussion}

\subsection{\textsc{\textmd{simba's}} Extragalactic Background Light}

The predicted \textsc{simba} EBL, integrated from $z = 0 \rightarrow 7$, is presented in \cref{fig:finalEBL}. The redshift cutoff at $z = 7$ is justified by the negligible contribution of higher redshift sources to the total signal, as demonstrated by the light cone wedges in \cref{fig:lcWedge}. The three spectral components optical/NIR, far IR, and radio are shown separately alongside the curated observational EBL data from Driver et al., which is itself assembled from a range of independent surveys and measurement techniques. The heterogeneous origin of these data points means that systematic offsets between individual datasets are not unexpected, and should be borne in mind when interpreting residuals between the model and the observations. Overall, the predicted spectrum broadly reproduces the shape and amplitude of the observed EBL across the sampled wavelength range particularly the two expected peaks in the optical and far IR discussed earlier. The optical/NIR component is moderately under predicted, the far IR peak is well reproduced, and the radio is systematically over predicted. Jackknife uncertainties are shown as shaded bands but are generally small enough to be visually unresolvable on the full plot.

\begin{figure*}
    \centering
    \includegraphics[width=\linewidth]{figures/ebl_final/ebl_full_jackknife.png}
    \caption{The \textsc{simba} predicted EBL integrated over $0 < z < 7$, with jackknife uncertainties shown as shaded bands. The three model components optical/NIR (blue), far IR dust continuum (red), and radio (green) are shown separately alongside the compiled observational data from Driver et al. (black points). The dust free component of the optical curve is also plotted (dark blue).}
    \label{fig:finalEBL}
\end{figure*}

To isolate the contributions from distinct cosmic epochs, the EBL was computed separately in three redshift bins: $0 < z < 1$, $1 < z < 3$, and $3 < z < 7$, each constructed from an independent light cone with its own jackknife error estimate (\cref{fig:finalBinnedEBL}). The sum of the three bins recovers the full EBL of \cref{fig:finalEBL}, and the decreasing amplitude with increasing redshift bin confirms that the $z < 3$ universe accounts for the dominant share of the total EBL across all three spectral components. Contributions from $3 < z < 7$ are of order $10^{-2}\,\mathrm{nW\,m^{-2}\,sr^{-1}}$, further supporting the decision to truncate the light cone at $z = 7$.

Several physically motivated trends are apparent in the binned decomposition. The optical/NIR emission is dominated by the $0 < z < 1$ bin as expected, the apparent brightness of stellar populations falls with distance, so the local universe contributes disproportionately to the integrated starlight background. The far IR component peaks in the $1 < z < 3$ bin, coinciding with the epoch of peak cosmic star formation activity. Photons travelling from these distances traverse greater dust column densities, with UV and optical radiation progressively reprocessed into the far IR, whose emission is then further redshifted to longer observed wavelengths. The radio emission is likewise most prominent at low redshift, where radio loud AGN are both abundant and subject to minimal cosmological dimming.

The individual redshift bin estimates carry noticeably larger jackknife uncertainties than the full redshift prediction. This is a consequence of cosmic variance, where each narrow bin subtends a smaller comoving volume, making the jackknife sub-volumes more sensitive to the stochastic distribution of bright sources. When integrated over the full redshift range, these fluctuations partially cancel across the larger ensemble of sub-volumes, reducing the relative variance in the total EBL estimate. This behaviour is consistent with the result that cosmic variance scales inversely with the square root of the survey volume \parencite{driverQuantifyingCosmicVariance2010}.

\begin{figure*}[htbp]
    \centering
    \includegraphics[width=0.9\linewidth]{figures/ebl_final/ebl_jackknife_bins.png}
    \caption{The \textsc{simba} predicted EBL decomposed into redshift bins $0 < z < 1$, $1 < z < 3$, and $3 < z < 7$, with jackknife uncertainties shown as shaded bands. The dominant contribution from $z < 3$ is evident across all three spectral components, with the $3 < z < 7$ bin contributing at the $10^{-2}\,\mathrm{nW\,m^{-2}\,sr^{-1}}$ level. The sum of the three bins recovers the total EBL.}
    \label{fig:finalBinnedEBL}
\end{figure*}

\subsection{The Optical/NIR Regime}

The \textsc{simba} simulations have demonstrated considerable success in reproducing key observational diagnostics relevant to the optical regime, including the stellar IMF, the galaxy stellar mass function, and broad star formation histories (SFHs) across cosmic time \parencite{daveSimbaCosmologicalSimulations2019}. Given this fidelity, the optical/NIR component of the predicted EBL was anticipated to be the most robustly constrained of the three regimes modelled here. While the overall spectral morphology and bin-by-bin behaviour are broadly consistent with expectation, the predicted amplitude systematically under predicts compared to the observational data.

Firstly, simulation snapshots were excluded from the integration due to missing \textsc{fsps} photometry in certain broadband filters, since the EBL was constructed by summing over all available snapshots, each excluded snapshot represents a gap in the cumulative light budget. This is likely to have a underestimating systematic effect on the total predicted flux, particularly if the missing snapshots correspond to epochs of elevated star formation activity.

Second, the reliability of the photometry used is contingent on the values computed and stored by \textsc{caesar} during the halo finding and galaxy cataloguing phase. The current pipeline adopts these pre-computed integrated galaxy properties rather than constructing a full SED per individual galaxy from first principles. This introduces a sensitivity to any systematic offsets or approximations present within \textsc{caesar}'s handling of stellar population properties. If the stored luminosities encode assumptions about mean stellar ages or metallicities that deviate from the true distribution within each galaxy, the integrated EBL will inherit those biases. A more robust approach would involve generating a bespoke SED for each simulated galaxy at each snapshot, applying stellar population synthesis (SPS) models conditioned on the individual metallicity and age distributions drawn directly from the SIMBA particle data. This would remove the dependence on \textsc{caesar}'s precomputed outputs entirely and would be the considered approach for future work.

The dust free optical component exhibits expected behaviour, with a systematically higher overall brightness compared to the dust attenuated case. This is physically consistent, as this excess emission is what is being absorbed by dust in the interstellar medium and subsequently reprocessed into the far infrared we plot. Furthermore, the intrinsic, dust free component is sufficiently high to closely trace the dust corrected observational data shown in \cref{fig:finalEBL}. This is possibly due to the dust attenuated photometry computed by \textsc{caesar} over correcting the emission, thereby causing the observed under prediction. This discrepancy is unlikely to stem from the underlying dust modelling within \textsc{simba}, which demonstrates robust agreement with theoretical expectations \cite{liDusttoGasDusttoMetalsRatio2019}. Notably, in the $1 < z < 3$ redshift bin, the dust free optical curve exhibits a significantly elevated flux at the low wavelength end closely matching the dust attenuated observational data. This behaviour further suggests that while the simulation accurately captures the intense young stellar populations expected at cosmic noon, the subsequent attenuation model by \textsc{caesar} is over estimating the dust extinction.

Despite these limitations, the redshift evolution remains physically consistent. In particular, the behaviour observed at high redshift ($3 < z < 7$) exhibits fluctuations at short wavelengths, which can be attributed to the redshifting of intrinsic ultraviolet emission from young, star-forming galaxies. At these redshifts, the observed frame wavelengths probe rest frame emission blueward of the Lyman break at $1216(1+z) \text{Å}$, where the contribution from massive, short-lived stars is significant. Additionally, the lower dust content in high-redshift galaxies reduces attenuation, allowing a larger fraction of ultraviolet radiation to escape and contribute to the observed SED. These effects naturally lead to enhanced variability and structure in the optical regime at high redshift.

% There could be a potential over or under correction in \textsc{simba} of the dust processes need sources for this not sure. should research into literature as to how good simba dust actually is and if this is even possible.
 
The dust attenuated peak is located at $\lambda \approx 1.19,\mu\mathrm{m}$ with an intensity of $\nu I_\nu = 7.31 \pm 0.52\mathrm{nWm^{-2}sr^{-1}}$, with a relative uncertainty of approximately $7.2\%$. Correspondingly, the non attenuated optical peak had a relative uncertainty of $2.4\%$. This small uncertainty indicates that the discrepancy between the simulated and observed flux is unlikely to be driven by sampling variance alone and it is not strongly sensitive to cosmic variance within the \textsc{simba} simulation volume. Instead pointing towards systematic effects in the modelling or data processing pipeline. 

\subsection{The Far Infrared Regime}

The far IR component of the predicted EBL is well reproduced overall, with the MBB prescription capturing the broad shape and amplitude of the observed far IR peak with reasonable accuracy.

An interesting feature of the far IR curve in \cref{fig:finalBinnedEBL} is its redshift dependence. The lowest redshift bins contribute more strongly at shorter wavelengths, while higher redshift bins dominate at longer wavelengths. This behaviour is expected because of two effects. Firstly, thermal dust emission is redshifted to larger $\lambda_\mathrm{obs}$ with increasing distance, and secondly dust star forming galaxies (DSFGs) dominate the early universe between $z=2\rightarrow6$ especially compared to the local universe, hence reprocessing a greater fraction of UV and optical radiation into the far IR \parencite{reuterCompleteRedshiftDistribution2020}

Although the single temperature MBB in \cref{eq:dustmbb} does reproduce the expected shape, representing each galaxy with a single effective temperature from \cref{eq:equivalentDustTemp} is a major simplification. Real galaxies host a distribution of grain temperatures set by different environments within the ISM, including diffuse dust, molecular clouds, and star-forming regions, where this temperature mixing is not well represented by a single value but rather a distribution. The peak and Wein side (blueward wavelengths) of the far IR SED are adequately matched, but contributions from cooler dust in the Rayleigh-Jeans tail may be underestimated due to this \parencite{caseyFarinfraredSpectralEnergy2012}. The predicted cosmic dust temperature of $44.7 \pm 7.7$\,K overestimates typical observed values, which generally peak between 35K and 40K \parencite{trainaA3COSMOSDustMass2024,ealesRiseFallDust2024}. However, the large standard deviation in this prediction is physically consistent with the scatter reported by \cite{liangDustTemperaturesHighredshift2019}, which is primarily driven by variations in star formation rates. Ultimately, we believe the overestimation is a systematic consequence of applying \cref{eq:equivalentDustTemp}. Because this relation was calibrated for dust temperatures at redshifts $z = 2-6$, it does not accurately capture the more quiescent conditions of the local universe. As discussed by \cite{parenteDriversCosmicDust2026}, local dust temperatures are intrinsically lower, therefore, extrapolating this high redshift prescription to local galaxies may artificially inflate the resulting temperature distribution.

The best fit value obtained from the fitting procedure shown in \cref{fig:tempDiagnostic} is $a_\mathrm{dust} = -0.0173 \pm 0.0006$, with a reduced $\chi^2 = 13.62$. The large reduced $\chi^2$ is interpreted as an indication of the limitations of the adopted model. The observational far IR EBL data from Driver et al. are compiled from multiple independent surveys, each subject to different systematic uncertainties. As a result, the quoted errors may underestimate the true scatter between data points, leading to an artificially inflated $\chi^2$. Additionally systematic residuals between the model and observations are expected due to the adopted MBB, as explained above. Despite this, the uncertainty on $a_\mathrm{dust}$ remains small, suggesting that the parameter is well constrained within the framework of the model.

% L_FIR energy budget

The current model neglects several physical processes that can become important in the far IR regime. In particular, the interaction between dust and the CMB is not explicitly included. At high redshift and long wavelengths, CMB heating can modify the effective dust temperature and suppress the contrast of dust emission relative to the background. Omitting it is expected to produce a minor overestimation on the long wavelength tail, particularly in the highest redshift bins \parencite{desertInterstellarDustEmission2022}.
Additionally, the adopted continuum MBB model does not account for spectral line emission, which can contribute to the total infrared output in sub millimetre wavelength ranges. While these effects are not expected to dominate the overall shape of the FIR peak, they may contribute to discrepancies in the detailed spectral behaviour \parencite{smailPotentialInfluenceFarinfrared2011}.

Another limitation arises in the mid infrared transition between the optical/NIR and far IR regimes dominated by emission from PAHs. As this component is not included in the MBB model, the gap is entirely expected within the scope of this work. A more complete treatment of the infrared regime would require full radiative transfer modelling, for example using tools such as \textsc{powderday}, which can self consistently model dust absorption, scattering, and emission across a broad wavelength range. Or some more recent approaches using machine learning techniques, such as in \cite{dobbelsPredictingGlobalFarinfrared2020} to infer the full global far IR SED of the universe.

Jackknife resampling gives a dominant far IR peak at $\lambda_\mathrm{obs} = 171.53\,\mu$m with intensity $\nu I_\nu = 12.10 \pm 0.27 \; \mathrm{nW\,m^{-2}\,sr^{-1}}$ which is a relative uncertainty of $2.2\%$. This is substantially lower than the $7.2\%$ obtained for the optical peak, consistent with the far IR EBL being dominated by a smaller number of bright, dusty galaxies whose contribution converges rapidly with simulation volume. The dominant uncertainties in the far IR are therefore systematic rather than statistical.

\subsection{The Radio Regime}

The radio component while matching the general shape, does systematically over predict relative to the observational data across the modelled wavelength range.

One plausible interpretation of the over prediction is that the observational radio datasets are incomplete rather than model error. Current radio surveys, even deep pointed observations, remain sensitivity limited and resolve only the brightest end of the source population \parencite{condonResolvingRadioSource2012, zottiRadioMillimeterContinuum2010}. An update of the faint population of star forming galaxies and low luminosity AGN below current detection thresholds would raise the observed EBL and reduce the apparent discrepancy. This is supported by \cite{tompkinsCosmicRadioBackground2023}, who explicitly report their measurements as empirical lower limits at all frequencies except 3\,GHz, the data points drawn from this study therefore represent a floor on the radio background rather than a complete census. \cite{murphyExcessHighfrequencyRadio2018} further demonstrate that even with deep stacking analyses at 1.4\,GHz, the recovered background remains a factor of $\sim$4.5 below the total sky brightness measured by ARCADE-2, underscoring the degree to which resolved source counts underestimate the true signal. This is further supported by the agreement of the radio curve with observations in the low redshift bins and the model picking up a lot more radio background at deeper redshift bins and so supporting a likely underestimation of radio population at high redshifts.

Furthermore, the \cite{thomasRadioGalaxyPopulation2021} calibration for radio emission is derived at a rest frame frequency of 1.4\,GHz, with spectral index corrections applied to scale to observed frequencies. \cite{tabatabaeiRadioSpectralEnergy2025} recently demonstrated that the radio SED shape of star forming galaxies is not fixed across redshift, variation is driven by changes in ISM conditions, interstellar magnetic field strengths, and cosmic ray propagation. A similar problem arises for the AGN emission where AGN cores typically exhibit flat spectra ($\alpha \sim 0$) due to synchrotron self absorption, while extended radio lobes produce steep spectra ($\alpha \sim 0.7-0.8$) from optically thin synchrotron emission \parencite{condonRadioEmissionNormal1992}. Treating all AGN emission with a single index conflates these physically distinct components, particularly at the frequencies where core and lobe emission exchange dominance. Allowing the spectral index to vary as a function of galaxy properties, or separating AGN into core and lobe components with dedicated SED templates, would provide a more accurate treatment. Applying locally calibrated relations for radio emission out to $z \sim 7$ therefore introduces a systematic uncertainty that is difficult to quantify without redshift dependent radio SED modelling.

The radio component exhibits noticeably larger jackknife uncertainties than either the optical or far IR in the $0 < z < 1$ bin. The relative uncertainty at the radio peak ($\lambda_\mathrm{obs} = 2997.92\,\mu$m) is $9.8\%$, compared to $7.2\%$ and $2.2\%$ for the optical and far-IR peaks respectively, and individual redshift bins reach higher values close to 37\%. This is a direct consequence of the radio EBL being dominated by AGN jet emission rather than star formation driven radiation. AGN are rare and spatially clustered, so their distribution across the jackknife sub-volumes is highly uneven, sub-volumes containing a bright radio AGN contribute disproportionately to the total signal, while adjacent sub-volumes do not. This spatial variance is inherently larger than for the far IR, where the signal is distributed across a much larger number of individually fainter star forming galaxies. The effect is most pronounced at low redshift, where the comoving volume sampled per snapshot is smallest and a single powerful AGN can dominate the local radio background.

\section{Conclusions}

In this work, a prediction of the EBL has been constructed using the \textsc{simba} cosmological hydrodynamical simulation, spanning the wavelength range from the optical/NIR to the radio regime. By combining galaxy catalogues with light cone techniques, a self-consistent estimate of the integrated CSED over the redshift range $0 \leq z \leq 7$ was obtained and compared to observational constraints.

The resulting EBL reproduces the expected bimodal structure of the cosmic energy distribution, with a primary peak arising from stellar emission in the optical/NIR and a secondary peak corresponding to dust reprocessing in the far infrared. The overall shape of the simulated spectrum is in good qualitative agreement with observations, indicating that the dominant physical processes governing photon production and energy redistribution in galaxies are broadly captured by the simulation framework.

%The comparison between the dust-attenuated and dust-free components indicates that the observed under-prediction is likely caused by the \textsc{caesar} model over-estimating dust extinction, as the intrinsic \textsc{simba} emission successfully matches observational data prior to attenuation.

The optical under prediction motivates the development of per galaxy SPS modelling conditioned on \textsc{simba} particle metallicity and age distributions, bypassing the pre computed \textsc{caesar} photometry. The radio over prediction, if not due to observational constraints, would suggest that \textsc{simba} AGN jet feedback prescription injects too much energy into the radio emitting electron population, with implications for the high redshift AGN luminosity function. The far infrared agreement, while partly by construction through the $a_{\mathrm{dust}}$ calibration, nonetheless demonstrates that \textsc{simba's} global dust energy budget is broadly consistent with observations, a non trivial result given that the dust-to-metal ratio evolves self-consistently across twelve billion years of cosmic time.

% Quantitatively, however, systematic discrepancies are present. The optical/NIR component is found to be under predicted, which is likely attributable to limitations in the use of precomputed photometry, missing simulation snapshots, and uncertainties in the modelling of dust attenuation. In contrast, the far infrared peak is well reproduced in both shape and normalisation, suggesting that the global dust energy budget in the simulation is reasonably consistent with observations, despite the simplified single temperature MBB approach. The radio component is over predicted relative to observational data, however this is likely believed to be due to incompleteness in the radio surveys. ogether, these results reinforce the interpretation that discrepancies in the optical regime are driven primarily by the treatment of dust attenuation rather than inaccuracies in the underlying stellar emission. 

The redshift analysis demonstrates that the majority of the EBL originates from $z < 3$, consistent with the peak of cosmic star formation activity. Optical emission is dominated by low redshift galaxies, while the far infrared receives significant contributions from intermediate redshifts due to the combined effects of dust reprocessing and cosmological redshift. High redshift galaxies contribute primarily through redshifted ultraviolet emission, producing the structured tail observed in the optical regime.

Statistical uncertainties, estimated via jackknife resampling, are found to be small for the total EBL, with relative errors of order a few percent at the peak wavelengths. This indicates that the dominant sources of discrepancy between the simulated and observed EBL are systematic in nature, rather than arising from sampling variance. In particular, simplified treatments of stellar emission, dust physics, and radio processes represent the primary limitations of the current approach.

Several avenues for future work emerge from this analysis. A more complete treatment of the infrared regime would require the inclusion of mid infrared emission, particularly from PAHs, which were beyond the scope of this study. Additionally, extending the analysis to higher energy regimes, including X-ray and gamma-ray emission would provide a more comprehensive picture of the total cosmic energy budget. Improvements to the modelling of radio emission, particularly through more physically motivated treatments of AGN jets and feedback, could also help resolve the observed discrepancies in long wavelength regimes.

Overall, this work demonstrates that cosmological hydrodynamical simulations such as \textsc{simba} are capable of reproducing the broad features of the extragalactic background light across multiple wavelength regimes. At the same time, the remaining discrepancies highlight the sensitivity of the cosmic SED to the detailed modelling of baryonic processes, particularly those involving dust and feedback. As both simulations and observations continue to improve, comparisons of this kind will provide increasingly powerful constraints on models of galaxy formation and the evolution of the universe.

\section{Acknowledgements}
The authors acknowledge helpful discussions with supervisor Romeel Davé. We thank Simon Driver for providing early access to the observational EBL data and respective article. We thank Robert Thompson for developing \textsc{caesar}, and the \textsc{yt} team for development and support of \textsc{yt}. We thank Eric Tittley for support with the Institute for Astronomy research computing cluster, Cuillin.
 
%%% References %%%
\renewcommand{\refname}{References}
\addcontentsline{toc}{section}{References}
\printbibliography

\end{document}

need to explain overall physics a lot more and explain things